\documentclass[lualatex,ja=standard,a4paper,12pt]{bxjsarticle}
\usepackage{amsmath, amssymb, amsthm}
\usepackage{geometry}
\usepackage{xcolor}
\usepackage{tcolorbox}
\usepackage{booktabs}
\usepackage{array}

\geometry{margin=25mm}

\tcbuselibrary{theorems, skins, breakable}

\newtcolorbox{defbox}[1]{
  colback=blue!5, colframe=blue!60!black,
  fonttitle=\bfseries, title={#1}, arc=3pt, breakable
}
\newtcolorbox{thmbox}[1]{
  colback=green!5, colframe=green!50!black,
  fonttitle=\bfseries, title={#1}, arc=3pt, breakable
}
\newtcolorbox{pointbox}{
  colback=orange!8, colframe=orange!70!black,
  fonttitle=\bfseries, title={ポイント}, arc=3pt, breakable
}
\newtcolorbox{exbox}[1]{
  colback=gray!6, colframe=gray!50,
  fonttitle=\bfseries, title={例：#1}, arc=3pt, breakable
}
\newtcolorbox{derivbox}{
  colback=purple!4, colframe=purple!50!black,
  fonttitle=\bfseries, title={導出}, arc=3pt, breakable
}

\title{{\large Hoff 第5章} \\[6pt]
       \textbf{5.5\; Prior Specification Based on Expectations} \\[4pt]
       {\normalsize ゼミ資料}}
\author{}
\date{}

\begin{document}
\maketitle
\thispagestyle{empty}

% ===========================
\section{この節のモチベーション}
% ===========================

前節までで正規モデルの共役事前分布として

\[
  \theta|\sigma^2 \sim \mathrm{Normal}(\mu_0,\, \sigma^2/\kappa_0),
  \qquad
  1/\sigma^2 \sim \mathrm{Gamma}(\nu_0/2,\, \nu_0\sigma_0^2/2)
\]

を使ってきた．ここでのハイパーパラメータ $(\mu_0, \kappa_0, \nu_0, \sigma_0^2)$ は
「事前サンプルサイズ」などと解釈されてきたが，実際に設定しようとすると難しい．

\medskip
\noindent
\textbf{自然な疑問：}「$Y$の期待値は $\mu_0$ くらい，分散は $\sigma_0^2$ くらい」
という形で事前情報を持っているとき，それをどうハイパーパラメータに変換するのか？

\medskip
\noindent
この節では，\textbf{観測量の期待値}（$E[Y]$ と $E[Y^2]$）を指定することで
ハイパーパラメータを自然に決める方法を導く．

% ===========================
\section{$p$次元指数型分布族}
% ===========================

\begin{defbox}{$p$次元指数型分布族}
密度が次の形で書けるモデルを $p$次元指数型分布族と呼ぶ：
\[
  p(y|\phi) = h(y)\,c(\phi)\exp\!\bigl\{\phi^\top t(y)\bigr\},
\]
ここで
\begin{itemize}
  \item $\phi = (\phi_1, \ldots, \phi_p)$：自然パラメータ
  \item $t(y) = (t_1(y), \ldots, t_p(y))$：十分統計量
  \item $c(\phi)$：正規化定数
\end{itemize}
\end{defbox}

\subsection*{正規モデルは2次元指数型分布族}

$Y \sim \mathrm{Normal}(\theta, \sigma^2)$ の密度を整理すると，

\[
  p(y|\theta,\sigma^2)
  = \frac{1}{\sqrt{2\pi\sigma^2}}\exp\!\left\{-\frac{(y-\theta)^2}{2\sigma^2}\right\}
  = \frac{1}{\sqrt{2\pi}}\exp\!\left\{\frac{\theta}{\sigma^2}y - \frac{1}{2\sigma^2}y^2
    - \frac{\theta^2}{2\sigma^2} - \frac{1}{2}\log\sigma^2\right\}
\]

なので，各成分は次のようになる：

\[
  t(y) = \begin{pmatrix}y \\ y^2\end{pmatrix},
  \qquad
  \phi = \begin{pmatrix}\theta/\sigma^2 \\ -1/(2\sigma^2)\end{pmatrix},
  \qquad
  c(\phi) = |\phi_2|^{1/2}\exp\!\left\{\frac{\phi_1^2}{4\phi_2}\right\}.
\]

\begin{pointbox}
十分統計量が $t(y) = (y,\, y^2)$ なので，$(\bar{y},\, s^2)$ を知ることは
$\bigl(\sum y_i,\, \sum y_i^2\bigr)$ を知ることと同値であり，推論に十分な情報が得られる．
\end{pointbox}

% ===========================
\section{自然パラメータ $\phi$ による共役事前分布}
% ===========================

第3章（1次元指数型分布族）と同様に，$p$次元でも共役事前分布は
$\phi$ の関数として次の形で書ける：

\[
  p(\phi \mid n_0, t_0)
  \propto c(\phi)^{n_0} \exp\!\bigl\{n_0\, t_0^\top \phi\bigr\},
\]

ここで $t_0 = (t_{01}, t_{02})$ は事前的な $Y$ の期待十分統計量，
すなわち

\[
  t_{01} = E[Y], \qquad t_{02} = E[Y^2].
\]

$n_0$ は「事前サンプルサイズ」に相当するハイパーパラメータである．

% ===========================
\section{$(\theta, \sigma^2)$ へのパラメータ変換}
% ===========================

自然パラメータ $\phi$ による共役事前分布を $(\theta, \sigma^2)$ で
書き直すと，大括弧内の指数部分は次の2つの因子に分解できる：

\begin{align*}
  p(\theta, \sigma^2 \mid n_0, t_0)
  &\propto
  (\sigma^2)^{-1/2}\exp\!\left\{-\frac{n_0(\theta - t_{01})^2}{2\sigma^2}\right\} \\
  &\quad \times\;
  (\sigma^2)^{-(n_0+5)/2}\exp\!\left\{-\frac{n_0(t_{02}-t_{01}^2)}{2\sigma^2}\right\}.
\end{align*}

\noindent
これはそれぞれ標準的な分布の核と一致する：

\medskip
\begin{center}
\renewcommand{\arraystretch}{1.6}
\begin{tabular}{lll}
\toprule
因子 & 分布の形 & 対応する分布 \\
\midrule
第1因子 & $\exp\!\left\{-\dfrac{n_0(\theta-t_{01})^2}{2\sigma^2}\right\}$ &
  $\theta|\sigma^2 \sim \mathrm{Normal}\!\left(t_{01},\, \dfrac{\sigma^2}{n_0}\right)$ \\[8pt]
第2因子 & $(\sigma^2)^{-(n_0+5)/2}\exp\!\left\{-\dfrac{n_0(t_{02}-t_{01}^2)}{2\sigma^2}\right\}$ &
  $\sigma^2 \sim \mathrm{Inv\text{-}Gamma}\!\left(\dfrac{n_0+3}{2},\; \dfrac{n_0(t_{02}-t_{01}^2)}{2}\right)$ \\
\bottomrule
\end{tabular}
\end{center}

% ===========================
\section{準備：反復期待値の法則と反復分散の法則}
% ===========================

以降の導出で繰り返し使う2つの法則を先にまとめる．

\begin{thmbox}{反復期待値の法則（Law of Total Expectation）}
確率変数 $X, Y$ に対して，
\[
  E[Y] = E\!\bigl[E[Y \mid X]\bigr].
\]
右辺の内側 $E[Y|X]$ は「$X$ を固定したときの $Y$ の平均」（$X$ の関数），
外側の $E[\,\cdot\,]$ はそれをさらに $X$ について平均したもの．
\end{thmbox}

\noindent\textbf{直感的な理解：}
クラス全体のテスト平均点 $E[Y]$ を求めたいとき，
「勉強した/しない」でグループに分けて各グループの平均を出し，
それをグループの割合で加重平均しても同じ答えが出る，ということ．

\medskip
\noindent\textbf{具体例：}テストの点数 $Y$，「勉強したか」$X$（1=した，0=しない）

\begin{center}
\renewcommand{\arraystretch}{1.3}
\begin{tabular}{lll}
\toprule
グループ & グループ内平均 & 割合 \\
\midrule
勉強した（$X=1$） & $E[Y|X=1] = 80$ & $P(X=1) = 0.7$ \\
しなかった（$X=0$） & $E[Y|X=0] = 50$ & $P(X=0) = 0.3$ \\
\bottomrule
\end{tabular}
\end{center}

\[
  E\!\bigl[E[Y|X]\bigr]
  = 0.7 \times 80 + 0.3 \times 50 = 71
  = E[Y].
\]

\begin{thmbox}{反復分散の法則（Law of Total Variance）}
\[
  \mathrm{Var}[Y]
  = \underbrace{E\!\bigl[\mathrm{Var}[Y \mid X]\bigr]}_{\text{グループ内分散の平均}}
  + \underbrace{\mathrm{Var}\!\bigl[E[Y \mid X]\bigr]}_{\text{グループ平均どうしの分散}}.
\]
\end{thmbox}

\noindent\textbf{直感的な理解：}
クラス全体の点数のバラつきは，
「各グループ内のバラつき（$X$ を固定したときの $Y$ の散らばり）」と
「グループ平均どうしのズレ（$E[Y|X=1]=80$ と $E[Y|X=0]=50$ の差）」
の2つが合わさったものである．

\medskip
\noindent\textbf{5.5節での使い方：}

\begin{center}
\renewcommand{\arraystretch}{1.5}
\begin{tabular}{ll}
\toprule
使う式 & 意味 \\
\midrule
$E[Y] = E[E[Y|\theta]] = E[\theta]$ & $\theta$ が不明でも，$\theta$ の平均が $Y$ の平均になる \\
$E[Y^2] = E[E[Y^2|\theta,\sigma^2]]$ & $Y^2$ の全体平均を，内側で $(\theta,\sigma^2)$ を固定して計算 \\
$\mathrm{Var}[\theta] = E[\mathrm{Var}[\theta|\sigma^2]]$ & $\theta$ の分散を $\sigma^2$ の不確かさで積分して求める \\
\bottomrule
\end{tabular}
\end{center}

% ===========================
\section{事前期待値から $t_{01}$，$t_{02}$ を決める}
% ===========================

\subsection*{「期待する」とはどういう意味か}

ハイパーパラメータを直接決めるのは難しい．$\mu_0, \sigma_0^2$ はモデル内部の量で，
観測できる $Y$ と直接対応していないからである．

\begin{thmbox}{発想の転換}
ハイパーパラメータを直接決める代わりに，
「\textbf{まだ見ていない新しい観測値 $Y$ に対して何を期待するか}」を
先に決め，それからハイパーパラメータを逆算する．
\end{thmbox}

具体的には：
\begin{itemize}
  \item \textbf{「$Y$ の平均はだいたい $\mu_0$ だろう」}\\
        （例：女性1人のIQは平均100程度 → $E[Y] = 100$）
  \item \textbf{「$Y$ は真の平均 $\theta$ の周りに $\sigma_0^2$ 程度バラつくだろう」}\\
        （例：個人のIQは母平均から $\pm 15$ 点程度 → $E[\mathrm{Var}[Y|\theta,\sigma^2]] = 225$）
\end{itemize}

この「期待する」は，$\theta$ と $\sigma^2$ の不確かさをすべて積分した
\textbf{周辺分布 $p(Y)$ での期待値}である．
つまり「$\theta$ の値を知らない状態で，$\theta \sim$ 事前分布，
$Y|\theta,\sigma^2 \sim \mathrm{Normal}(\theta,\sigma^2)$ という
2段階の確率の流れを通じた期待値」を意味する．

\subsection*{正規モデルの条件付き期待値・分散}

導出で繰り返し使う事実をまとめる．
正規モデル $Y|\theta,\sigma^2 \sim \mathrm{Normal}(\theta,\sigma^2)$，
事前分布 $\theta|\sigma^2 \sim \mathrm{Normal}(\mu_0, \sigma^2/n_0)$ から直接読める：

\begin{center}
\renewcommand{\arraystretch}{1.5}
\begin{tabular}{lll}
\toprule
条件付き量 & 値 & 根拠 \\
\midrule
$E[Y|\theta,\sigma^2]$ & $\theta$ & 正規分布の平均はパラメータ $\theta$ \\
$\mathrm{Var}[Y|\theta,\sigma^2]$ & $\sigma^2$ & 正規分布の分散はパラメータ $\sigma^2$ \\
$E[\theta|\sigma^2]$ & $\mu_0$ & 正規分布の平均はパラメータ $\mu_0$ \\
$\mathrm{Var}[\theta|\sigma^2]$ & $\sigma^2/n_0$ & 正規分布の分散はパラメータ $\sigma^2/n_0$ \\
\bottomrule
\end{tabular}
\end{center}

\noindent これらは「正規分布だから条件付き期待値・分散はパラメータそのもの」という
単純な事実であり，法則ではなくモデルの定義から直接来る．

\subsection*{$t_{01}$ の決め方}

「$Y$ の平均が $\mu_0$」すなわち $E[Y] = \mu_0$ を期待するとき：

\[
  E[Y]
  \underset{\text{反復期待値の法則}}{=} E\!\bigl[E[Y|\theta]\bigr]
  \underset{E[Y|\theta]=\theta}{=} E[\theta]
  = \mu_0.
\]

ゆえに $t_{01} = \mu_0$．

\subsection*{$t_{02}$ の決め方}

「$Y$ の $\theta$ 周りのバラつきが $\sigma_0^2$」すなわち
$E[\mathrm{Var}[Y|\theta,\sigma^2]] = \sigma_0^2$ を期待するとき：

\[
  E\!\bigl[\mathrm{Var}[Y|\theta,\sigma^2]\bigr]
  \underset{\mathrm{Var}[Y|\theta,\sigma^2]=\sigma^2}{=} E[\sigma^2] = \sigma_0^2.
\]

これは法則ではなく，正規分布の分散がパラメータ $\sigma^2$ そのものであることを使った代入である．

\medskip
$t_{02} = E[Y^2]$ を求める．まず $E[Y^2|\theta,\sigma^2]$ を整理する：
\[
  E[Y^2|\theta,\sigma^2]
  = \underbrace{\mathrm{Var}[Y|\theta,\sigma^2]}_{\sigma^2}
  + \underbrace{(E[Y|\theta,\sigma^2])^2}_{\theta^2}
  = \sigma^2 + \theta^2.
\]

\begin{derivbox}
\begin{align*}
  t_{02} = E[Y^2]
  &\underset{\text{反復期待値の法則}}{=} E\!\bigl[E[Y^2|\theta,\sigma^2]\bigr]
   \underset{\text{上の結果を代入}}{=} E[\sigma^2 + \theta^2]
   = \underbrace{E[\sigma^2]}_{\sigma_0^2} + E[\theta^2].
\end{align*}

$E[\theta^2] = \mathrm{Var}[\theta] + \mu_0^2$ であり，$\mathrm{Var}[\theta]$ は
反復分散の法則と $E[\theta|\sigma^2]=\mu_0$，$\mathrm{Var}[\theta|\sigma^2]=\sigma^2/n_0$ から：
\[
  \mathrm{Var}[\theta]
  = E\!\bigl[\underbrace{\mathrm{Var}[\theta|\sigma^2]}_{\sigma^2/n_0}\bigr]
  + \mathrm{Var}\!\bigl[\underbrace{E[\theta|\sigma^2]}_{\mu_0=\text{定数}}\bigr]
  = \frac{\sigma_0^2}{n_0} + 0
  = \frac{\sigma_0^2}{n_0}.
\]

したがって：
\[
  t_{02} = \sigma_0^2 + \frac{\sigma_0^2}{n_0} + \mu_0^2
         = \sigma_0^2\cdot\frac{n_0+1}{n_0} + \mu_0^2,
  \qquad
  n_0(t_{02} - t_{01}^2) = (n_0+1)\sigma_0^2.
\]
\end{derivbox}

\begin{pointbox}
どちらの導出も同じ構造：
\begin{enumerate}
  \item \textbf{反復期待値（分散）の法則}で $(\theta,\sigma^2)$ についての積分を分割する
  \item \textbf{正規分布の性質}（条件付き期待値・分散 = パラメータそのもの）を代入して計算する
\end{enumerate}
「期待する」という主観的な量（$\mu_0$，$\sigma_0^2$）は，
この2ステップを経てハイパーパラメータ $t_{01}, t_{02}$ に翻訳される．
\end{pointbox}

% ===========================
\section{得られる事前分布}
% ===========================

以上をまとめると，事前期待値 $\mu_0 = E[Y]$，$\sigma_0^2 = E[\mathrm{Var}[Y|\theta,\sigma^2]]$
を指定したとき，対応する共役事前分布は

\begin{thmbox}{事前分布（期待値ベースのパラメータ化）}
\[
  \theta\,|\,\sigma^2 \sim \mathrm{Normal}\!\left(\mu_0,\; \frac{\sigma^2}{n_0}\right)
\]
\[
  \sigma^2 \sim \mathrm{Inv\text{-}Gamma}\!\left(\frac{n_0+3}{2},\; \frac{(n_0+1)\sigma_0^2}{2}\right)
\]
\end{thmbox}

\noindent
共役事前分布の一般形は $\kappa_0$（$\theta$ の事前精度）と $\nu_0$（$\sigma^2$ の自由度）を
\textbf{独立に}設定できるが，5.5節のアプローチはこれらを $n_0$ ひとつに縛る：

\begin{center}
\renewcommand{\arraystretch}{1.4}
\begin{tabular}{lll}
\toprule
 & 一般の共役事前分布 & 5.5節のパラメータ化 \\
\midrule
$\theta$ の事前精度 & $\kappa_0$（自由） & $n_0$ \\
$\sigma^2$ の自由度 & $\nu_0$（自由） & $n_0 + 3$ \\
\bottomrule
\end{tabular}
\end{center}

\noindent 暗黙に $\nu_0 = \kappa_0 + 3$ という制約が入るため，
自由度は $(\kappa_0, \nu_0)$ の2つから $n_0$ の1つに減る．

\begin{pointbox}
\textbf{トレードオフ：}
\begin{itemize}
  \item \textbf{利点}：$n_0$「事前サンプルサイズ」として直感的に解釈できる
  \item \textbf{欠点}：「$\theta$ への確信は強いが $\sigma^2$ への確信は弱い」
        という非対称な設定ができない
\end{itemize}
解釈しやすさのために柔軟性を1自由度分だけ犠牲にした設計である．
\end{pointbox}

\subsection*{弱い事前分布の場合：$n_0 = 1$}

事前情報が少ない場合は $n_0 = 1$ とする：

\[
  \theta|\sigma^2 \sim \mathrm{Normal}(\mu_0,\; \sigma^2),
  \qquad
  1/\sigma^2 \sim \mathrm{Gamma}\!\left(2,\; \sigma_0^2\right).
\]

この設定のもとで，事前期待値が正しく再現されていることを確認できる：

\[
  E[Y] = E[E[Y|\theta]] = E[\theta] = \mu_0 \quad \checkmark
\]
\[
  E[\mathrm{Var}[Y|\theta,\sigma^2]] = E[\sigma^2] = \sigma_0^2 \quad \checkmark
\]

% ===========================
\section{事後分布}
% ===========================

$n_0 = 1$ の事前分布のもとで $n$個の i.i.d.\ サンプルを得たとき，
事後分布は次のようになる：

\begin{thmbox}{事後分布（$n_0=1$ の場合）}
\[
  \{\theta\,|\,\sigma^2, y_1,\ldots,y_n\}
  \sim \mathrm{Normal}\!\left(
    \frac{\mu_0/\sigma^2 + n\bar{y}}{1/\sigma^2 + n},\;\;
    \frac{\sigma^2}{n+1}
  \right)
\]
\[
  \{\sigma^2\,|\,y_1,\ldots,y_n\}
  \sim \mathrm{Inv\text{-}Gamma}\!\left(
    2 + \frac{n}{2},\;\;
    \sigma_0^2 + (n-1)s^2 + \frac{n}{n+1}(\bar{y}-\mu_0)^2
  \right)
\]
\end{thmbox}

\subsection*{事後平均の解釈}

$\theta$ の事後平均は

\[
  \frac{\mu_0/\sigma^2 + n\bar{y}}{1/\sigma^2 + n}
  = \underbrace{\frac{1/\sigma^2}{1/\sigma^2 + n}}_{\text{事前の重み}}\mu_0
  + \underbrace{\frac{n}{1/\sigma^2 + n}}_{\text{データの重み}}\bar{y}.
\]

第5.3節の公式と同じく，事前平均とデータの加重平均になっている．
$n \to \infty$ でデータの重みが1に収束し，事前情報の影響が消える．

\subsection*{$\sigma^2$ の事後スケールパラメータの解釈}

スケールパラメータ $\sigma_0^2 + (n-1)s^2 + \dfrac{n}{n+1}(\bar{y}-\mu_0)^2$ の各項：

\medskip
\begin{center}
\renewcommand{\arraystretch}{1.5}
\begin{tabular}{ll}
\toprule
項 & 意味 \\
\midrule
$\sigma_0^2$ & 事前のスケール（事前サンプルサイズ相当の分散） \\
$(n-1)s^2$ & データの変動（サンプル内の散らばり） \\
$\dfrac{n}{n+1}(\bar{y}-\mu_0)^2$ & 事前平均とデータ平均のズレ \\
\bottomrule
\end{tabular}
\end{center}

\medskip
3つ目の項は，事前平均とデータの乖離が大きいほど $\sigma^2$ が大きい方向に
更新されることを意味する（「想定外にズレているなら分散も大きいはず」という情報）．

% ===========================
\section{まとめ}
% ===========================

\begin{thmbox}{5.5節のまとめ}
\begin{enumerate}
  \item 正規モデルは2次元指数型分布族であり，十分統計量は $t(y) = (y, y^2)$，
        すなわち $(\bar{y}, s^2)$ に相当する．

  \item 共役事前分布を自然パラメータ $\phi$ で書くと，
        ハイパーパラメータ $t_0$ が観測量 $Y$ の\textbf{事前期待値}（$E[Y]$，$E[Y^2]$）
        として解釈できる．

  \item 「$Y$の期待値は $\mu_0$，分散の期待値は $\sigma_0^2$」という
        \textbf{解釈しやすい量}を指定するだけで，$(t_{01}, t_{02})$ が自動的に決まる．

  \item 事前サンプルサイズ $n_0$ だけが調整パラメータとして残り，
        「事前情報の強さ」を表す．

  \item 得られた事前・事後分布は第5.3節の共役モデルと同じ構造を持ち，
        事後平均は事前平均とデータ平均の加重平均になる．
\end{enumerate}
\end{thmbox}

\end{document}
